\section{Algorithm}\label{sec:alg}
To train group DRO models efficiently, we introduce an online optimization algorithm with convergence guarantees.
Prior work on group DRO has either used batch optimization algorithms, which do not scale to large datasets, or stochastic optimization algorithms without convergence guarantees.

In the convex and batch case, there is a rich literature on distributionally robust optimization which treats the problem as a standard convex conic program \citep{bental2013robust, duchi2016, bertsimas2018data, lam2015quantifying}.
For general non-convex DRO problems, two types of stochastic optimization methods have been proposed:
(i) stochastic gradient descent (SGD) on the Lagrangian dual of the objective \citep{duchi2018learning,hashimoto2018repeated}, and
(ii) direct minimax optimization \citep{namkoong2016stochastic}.
The first approach fails for group DRO because the gradient of the dual objective is difficult to estimate in a stochastic and unbiased manner.%
\footnote{
The dual optimization problem for group DRO is $\min_{\theta,\beta}\frac{1}{\alpha}\E_g[\max(0,\E_{x,y\sim\hP_g}\left[\ell(\theta; (x,y))\mid g\right]-\beta)] + \beta$ for constant $\alpha$. The $\max$ over \emph{expected} loss makes it difficult to obtain an unbiased, stochastic gradient estimate.
}
An algorithm of the second type has been proposed for group DRO \citep{oren2019drolm},
but this work does not provide convergence guarantees, and we observed instability in practice under some settings.

Recall that we aim to solve the optimization problem \refeqn{htdro}, which can be rewritten as
\begin{align}
  \min_{\theta \in \Theta} \sup_{q \in \Delta_m} \ \sum_{g=1}^m q_g \E_{(x, y) \sim P_g}[\ell(\theta; (x, y))].
\end{align}
Extending existing minimax algorithms for DRO \citep{namkoong2016stochastic,oren2019drolm},
we interleave gradient-based updates on $\theta$ and $\padv$.
Intuitively, we maintain a distribution $\padv$ over groups, with high masses on high-loss groups,
and update on each example proportionally to the mass on its group.
Concretely, we interleave SGD on $\theta$ and exponentiated gradient ascent on $\padv$
(\refalg{opt}). (In practice, we use minibatches and a momentum term for $\theta$; see \refapp{model_details} for details.)
The key improvement from the existing group DRO algorithm \citep{oren2019drolm} is that $\padv$ is updated using gradients instead of picking the group with worst average loss at each iteration,
which is important for stability and obtaining convergence guarantees.
The run time of the algorithm is similar to that of SGD for a given number of epochs (less than a 5\% difference), as run time is dominated by the computation of the loss and its gradient.

\begin{algorithm}[H]
\label{alg:opt}
 \DontPrintSemicolon
 \KwIn{Step sizes $\eta_q, \eta_\theta; P_g$ for each $g\in \gset$}
 Initialize $\theta\iter{0}$ and $\padv\iter{0}$\;
 \For{$t=1,\dots,T$}{
  $ g \sim \text{Uniform}(1,\dots,\ngroup)$ \tcp*{Choose a group $g$ at random}
  $ x,y \sim P_g$ \tcp*{Sample $x, y$ from group $g$}
  $ \padv^\prime \gets \padv\iter{t-1}$;\
  $ \padv^\prime_g \gets \padv^\prime_g\exp({\eta_q\ell(\theta\iter{t-1}; (x,y))})$ \tcp*{Update weights for group $g$}
  $ \padv\iter{t} \gets \padv^\prime / \sum_{g^\prime} \padv^\prime_{g^\prime}$ \tcp*{Renormalize $q$}
  $\theta\iter{t} \gets \theta\iter{t-1} - \eta_\theta \padv_g^{(t)}\nabla \ell(\theta\iter{t-1}; (x,y))$ \tcp*{Use $q$ to update $\theta$}
 }
  \caption{Online optimization algorithm for group DRO}
\end{algorithm}

We analyze the convergence rate by studying the error $\varepsilon_T$ of the average iterate $\bar{\theta}\iter{1:T}$:
\begin{align}
    \varepsilon_T
    &= \max_{\padv\in\Delta_\ngroup}L\bigl(\bar{\theta}\iter{1:T}, \padv\bigr)
       - \min_{\theta\in\Theta}\max_{\padv\in\Delta_\ngroup} L\bigl(\theta, \padv\bigr),
\end{align}
where $L(\theta, \padv) \eqdef \sum_{g=1}^\ngroup \padv_g\E_{(x, y) \sim P_g}[\ell(\theta; (x, y))]$ is the expected worst-case loss.
Applying results from \citet{nemirovski2009robust}, we can show that
\refalg{opt} has a standard convergence rate of $O\bigl(1 / \sqrt{T} \bigr)$ in the convex setting (proof in \refsec{proof_opt}):
\begin{proposition}\label{thm:opt}
  Suppose that the loss $\ell(\cdot; (x, y))$ is non-negative, convex, $\Bgrad$-Lipschitz continuous, and bounded by $\Bloss$ for all $(x, y)$ in $\sX \times \sY$, and $\|\theta\|_2\le \Btheta$ for all $\theta\in\Theta$ with convex $\Theta \subseteq \R^d$.
  Then, the average iterate of \refalg{opt} achieves an expected error at the rate
  \begin{align}
    \E[\varepsilon_T]
      \le 2\ngroup\sqrt{\frac{10(\Btheta^2\Bgrad^2 + \Bloss^2\log\ngroup)}{T}}.
  \end{align}
  where the expectation is taken over the randomness of the algorithm.
\end{proposition}
